% ===========================================================================
\section{Inversion, and a hard limit}\label{sec:why}
\warmth{0}\quad\whisper{Orthogonality gives the signal back. And no signal is sharp in both domains at once.}

\begin{strand}
Two facts anchor the theory. The transform is \emph{invertible} (the pure tones are orthonormal,
so coordinates determine the vector), and it obeys a \emph{trade-off} (a signal and its transform
cannot both be concentrated). The first is why Fourier is useful; the second is why signal
processing is interesting.
\end{strand}

\block{Getting the signal back}

\begin{theorem}[inversion]\label{thm:inv}
If $\widehat f$ is nice enough, then
\[
  f(x)\;=\;\int_{-\infty}^{\infty}\widehat f(\xi)\,e^{2\pi i\xi x}\dd\xi .
\]
That is: $f$ is the superposition of its tones, each weighted by its coordinate $\widehat f(\xi)$.
\end{theorem}
\begin{proof}
Skeleton: the tones are orthonormal in the sense $\int e^{2\pi i\xi x}\,\overline{e^{2\pi i\xi' x}}\dd x
=\delta(\xi-\xi')$. \TODO{substitute $\widehat f(\xi)=\int f(y)e^{-2\pi i\xi y}\dd y$ into the right side and collapse the $\xi$-integral with the orthogonality relation}
\end{proof}

\recall{So $\F$ is a unitary change of basis on $L^2(\R)$: it preserves lengths (Parseval) and is undone by $\F^{-1}$, which is the same formula with a $+$ sign in the exponent.}

\block{The uncertainty principle}

\begin{theorem}[Heisenberg uncertainty]\label{thm:unc}
For a unit-energy signal, let $(\Delta x)^2$ and $(\Delta\xi)^2$ be the variances of $\abs{f}^2$ and
$\abs{\widehat f}^2$. Then
\[
  \Delta x\,\cdot\,\Delta\xi\;\ge\;\frac{1}{4\pi},
\]
with \keyword{equality exactly for Gaussians}. You cannot be sharply localized in space and in
frequency at once.
\end{theorem}

\begin{strand}
This is not a defect of the transform; it is a law. A pure tone (one frequency) is spread over all
space; a spike (one instant) needs all frequencies. The Gaussian sits exactly at the boundary,
equally (and minimally) spread in both, which is why it is the transform's fixed point and the
shape of every ``optimal'' window, wave packet, and kernel.
\end{strand}

\begin{yourturn}
Scale a Gaussian: let $f_a(x)=e^{-\pi a x^2}$.
\begin{itemize}[leftmargin=1.3em,itemsep=1pt]
  \item Its transform is $\widehat f_a(\xi)=\tfrac{1}{\sqrt a}\,e^{-\pi\xi^2/a}$. As $a\to\infty$ the
        signal gets \fillin[2cm] in $x$ and \fillin[2cm] in $\xi$.
  \item So the product of the two widths stays \fillin[2.5cm] — the uncertainty bound is tight here.
\end{itemize}
\recall{Narrower in $x$, wider in $\xi$ (and vice versa); the width-product is constant. Squeeze one domain, the other balloons.}
\end{yourturn}

\begin{strand}
The same inequality, read in quantum mechanics with $x=$ position and $\xi=$ momentum, is
Heisenberg's. Position and momentum are a Fourier pair; the ``principle'' is a theorem about $\F$.
\end{strand}
